\documentclass[aspectratio=169,10pt,fontset=fandol]{ctexbeamer}
\usepackage{amsmath,amssymb,mathtools,bm,mathrsfs}
\usepackage{graphicx,booktabs,array,tabularx,multirow}
\usepackage{tikz,pgfplots,xparse,xstring,mdframed}
\usetikzlibrary{arrows.meta,calc,angles,positioning,shapes.geometric,graphs}
\pgfplotsset{compat=1.18}

\setbeamerfont{normal text}{size=\normalsize}
\setbeamerfont{title}{size=\large}
\setbeamerfont{subtitle}{size=\normalsize}
\setbeamerfont{frametitle}{size=\large}
\setbeamerfont{framesubtitle}{size=\small}
\setbeamerfont{block title}{size=\normalsize}
\setbeamerfont{block body}{size=\normalsize}
\setbeamerfont{section in toc}{size=\normalsize}
\setbeamerfont{section in head/foot}{size=\footnotesize}
\setbeamerfont{author in head/foot}{size=\footnotesize}
\setbeamerfont{title in head/foot}{size=\footnotesize}
\setbeamerfont{date in head/foot}{size=\footnotesize}
\setbeamertemplate{navigation symbols}{}
\setbeamertemplate{itemize item}{\normalsize\textbullet}
\setbeamertemplate{itemize subitem}{\small\textbullet}
\setbeamercolor{structure}{fg=blue!65!black}
\setbeamercolor{block title}{fg=blue!65!black,bg=blue!8}
\setbeamercolor{block body}{fg=black,bg=blue!3}
\setbeamertemplate{footline}{\leavevmode\hbox to\paperwidth{\hfill{\footnotesize\insertframenumber}\hspace*{1em}}\vskip3pt}

% ctex fontset=fandol; system fonts are not redistributed.

\newenvironment{proofdetail}{\par\noindent\textbf{证明：}\quad}{\par}
\renewenvironment{solution}{\par\noindent\textbf{解：}\quad}{\par}
\title{删除与连贯分页.tex}
\begin{document}
\section{保留计算}
% FRAME-ID F001
\typeout{MATH-AGENT-FRAME-BEGIN F001}
\begin{frame}[t,fragile]{目录}
\normalsize
{\color{blue}\normalsize 1. 保留计算\par}\vspace{0.22em}
{\color{gray!40}\normalsize 2. 两个方面\par}\vspace{0.22em}
\end{frame}
\typeout{MATH-AGENT-FRAME-END F001}
% FRAME-ID F002
\typeout{MATH-AGENT-FRAME-BEGIN F002}
\begin{frame}[t,fragile]{保留计算}
\normalsize
\begin{block}{例题}
{\normalsize
% SOURCE-BEGIN U0002 0 8
求 $1+1$。
% SOURCE-END U0002
}
\end{block}
\pause
{\normalsize
\textbf{解：}
% SOURCE-BEGIN U0003 0 8
答案为 $2$。
% SOURCE-END U0003
}
\end{frame}
\typeout{MATH-AGENT-FRAME-END F002}
\section{两个方面}
% FRAME-ID F003
\typeout{MATH-AGENT-FRAME-BEGIN F003}
\begin{frame}[t,fragile]{目录}
\normalsize
{\color{black}\normalsize 1. 保留计算\par}\vspace{0.22em}
{\color{blue}\normalsize 2. 两个方面\par}\vspace{0.22em}
\end{frame}
\typeout{MATH-AGENT-FRAME-END F003}
% FRAME-ID F004
\typeout{MATH-AGENT-FRAME-BEGIN F004}
\begin{frame}[t,fragile]{两个方面}
\normalsize
{\normalsize
% SOURCE-BEGIN U0004 0 20
第一，平均变化率描述一段区间。 


% SOURCE-END U0004
}
{\normalsize
% SOURCE-BEGIN U0005 0 17
第二，瞬时变化率描述一个时刻。

% SOURCE-END U0005
}
\end{frame}
\typeout{MATH-AGENT-FRAME-END F004}
\end{document}
